Citations provide the basis for trusting scientific claims; when they are invalid or fabricated, this trust collapses.
With the advent of Large Language Models (LLMs), this risk has intensified: LLMs are increasingly used for academic writing, yet their tendency to fabricate citations (``ghost citations'') poses a systemic threat to citation validity.
To quantify this threat and inform mitigation, we develop \citeb, an open-source framework for large-scale citation verification, and conduct a comprehensive study of citation validity in the LLM era through three complementary experiments.
We benchmark 13 state-of-the-art LLMs on citation generation across 40 research domains, finding that all models hallucinate citations with significant variation across domains.
Moreover, we analyze 2.2 million citations from 56,381 papers at top-tier AI/ML and Security venues (2020--2025), finding that 1.07\% of papers contain invalid citations, with rates increasing sharply in 2025.
Furthermore, we survey 97 researchers, revealing a critical verification gap between stated diligence and actual practices.
Our findings indicate that unreliable AI tools, combined with inadequate human verification, facilitate the penetration of fabricated citations into the scientific record.
We discuss implications and propose interventions for researchers, venues, and tool developers.
